\section{Related Works}
\label{sec:related}

% Four strands, all present: RV modeling; ML in prediction;
% regularization geometry; retransformation. The protocol-specific
% literature discussion sits with the estimator in Section 4.3; this
% section orients the paper.

\paragraph{Realized-volatility modeling.}
The paper's target and backbone descend from the realized-measure
literature: realized variance as a sum of squared intraday returns
\citep{AndersenBollerslev1997,BarndorffNielsenShephard2004}, its
approximate long memory documented by
\citet{AndersenBollerslevDieboldLabys2003}, and the heterogeneous
autoregressive model of \citet{Corsi2009}, whose daily/weekly/monthly
averages we adapt to a base-2 intraday ladder. Refinements of the HAR
backbone are numerous: measurement-error-dependent coefficients
\citep{BollerslevPattonQuaedvlieg2016}, state-dependent time variation
\citep{BekiermanManner2018}, decoupled short- and long-run behavior
\citep{BennedsenLundePakkanen2022}, joint models of returns and realized
measures \citep{HansenHuangShek2012}, and multiple-indicator intraday
structures \citep{EngleGallo2006}. We keep the backbone deliberately
plain --- per-bar least squares on the lag ladder --- because the paper's
questions concern what can be added to it, and its strength makes every
addition a stringent test: on this panel the HAR backbone's coefficients are
$94$--$97\%$ directionally stable over two decades, so a predictor that
survives next to it carries information the ladder does not.

\paragraph{Machine learning for volatility and return prediction.}
Tree ensembles are strong off-the-shelf estimators in the recent
machine-learning literature on asset pricing and volatility
\citep{gu2020empirical,christensen2023machine}, typically in
gradient-boosted form \citep{Ke2017,Chen2016}. Interpretability
machinery --- functional ANOVA, additive and pairwise additive models
\citep{Lou2013}, and Shapley-based attribution
\citep{lundberg2017unified,lundberg2018consistent} --- is usually
applied \emph{post hoc}. This paper uses the trees as a specification
search. The structure they recover --- sparse pairwise products gated
on the session clock --- is written back into the linear design as
explicit columns, and the distilled model is estimated in the same
rolling least-squares path as every other arm.

\paragraph{Sparse selection versus dense shrinkage.}
The linear estimators are standard families --- ridge
\citep{HoerlKennard1970}, lasso \citep{Tibshirani1996}, elastic net
\citep{ZouHastie2005}, with the least-angle-regression path
\citep{Efron2004} behind our exact online lasso solve --- but the
question they are put to has a now-substantial literature: whether
economic signal is concentrated in a few predictors or spread thinly
across many. \citet{GiannoneLenzaPrimiceri2021} compare sparse and
dense representations across macro, micro, and finance prediction
problems and find that the posterior rarely concentrates on a sparse
model --- sparsity is largely an illusion, and dense methods (ridge,
factor models) predict best. \citet{KozakNagelSantosh2020} reach the
parallel conclusion for the cross-section of expected returns, where a
dense, heavily shrunk stochastic discount factor dominates sparse
characteristic selections. \citet{ShenXiu2025} supply the estimation
theory: in the weak-signal regime the zero predictor is asymptotically
Bayes-optimal, ridge can improve on it while lasso cannot for any
penalty, and the ranking extends to nonlinear methods (dense random
forests over sparse boosted trees, $\ell_2$- over $\ell_1$-penalized
networks); lasso's failure reflects signal \emph{weakness}, not
necessarily the absence of sparsity. The rare-and-weak detection
literature \citep{DonohoJin2004} formalizes the underlying phase
structure, in which a dense regime of individually undetectable
effects is separated from the sparse regime by a sharp boundary. Our
head-to-head (Section~\ref{sec:dense_weak}) is the
intraday-volatility instance of this phenomenon, measured under a
stricter protocol than the literature's cross-validation standard:
every penalty is selected causally on a short forward tail, and the
comparison is scored on a consistent loss. That protocol isolates a
mechanism the asymptotic results abstract from: on short validation
tails the price of a hyperparameter mistake is asymmetric across
families --- benign for ridge, whose loss is flat in the penalty near
the optimum, destructive for selectors, whose penalty is partly a
selection dial --- which is why the paper's final model carries fixed
penalties rather than tuned ones.

\paragraph{Retransformation.}
A mean model estimated on a log or square-root scale and scored on
the original scale is not inverted by the algebraic inverse: by
Jensen's inequality the na\"ive back-transform is biased for the
original-scale mean. The standard nonparametric correction is
\citet{Duan1983}'s smearing estimator. \citet{Manning1998} and
\citet{ManningMullahy2001} place that correction inside the forecast:
under heteroscedasticity a single smearing factor is itself biased, so
the retransformation is a modeling choice rather than a reporting
footnote. This paper fits on a diurnally-adjusted square-root scale
and scores QLIKE on raw variance, so the same gap applies. The
back-transform is treated as part of the evaluation protocol --- a
calibrated, strictly causal second-moment correction that replaces the
retrospective in-window Duan scalar --- and the comparison is reported
in Section~\ref{sec:smearing_contract}.